% questions.tex —— 2014 年·山东卷·文科 · 题目内容片段（无答案）
% 由 tools/generate_projects.py 依来源总账自 2014/2014_数学真题_山东文_试题.md 生成，勿手改（改源 MD 后重跑）


\begin{gknotice}
\textbf{来源说明}：本题面依据固定源 PDF 的两页 A3 题面逐题转录，原卷印刷页码为 909—910。第 909 面为选择题 1—10、填空题 11—15 与解答题 16—17；第 910 面为解答题 18—21。全卷三处节头依次为「一、选择题」「二、填空题」「三、解答题」；题号与索引一致。

\textbf{解析归属}：固定源不附答案或解析；配套答案与解析均由本库按本题面独立复算。

\textbf{图形说明}：第 6 题的对数函数图象、第 8 题的频率分布直方图、第 11 题的程序框图、第 18 题的四棱锥由项目内 TikZ 矢量源重绘；其余各题原页无图。
\end{gknotice}

\gkpart{一、选择题}

\begin{question}
\gkti{1}已知 $a,b\in\mathbf R$，$\mathrm i$ 是虚数单位，若 $a+\mathrm i=2-b\mathrm i$，则 $(a+b\mathrm i)^2=$\mbox{（　　）}

（A）$3-4\mathrm i$；（B）$3+4\mathrm i$；（C）$4-3\mathrm i$；（D）$4+3\mathrm i$。

\end{question}
\begin{question}
\gkti{2}设集合 $A=\{x\mid x^2-2x<0\}$，$B=\{x\mid1\leqslant x\leqslant4\}$，则 $A\cap B=$\mbox{（　　）}

（A）$(0,2]$；（B）$(1,2)$；（C）$[1,2)$；（D）$(1,4)$。

\end{question}
\begin{question}
\gkti{3}函数 $f(x)=\dfrac1{\sqrt{\log_2x-1}}$ 的定义域为\mbox{（　　）}

（A）$(0,2)$；（B）$(0,2]$；（C）$(2,+\infty)$；（D）$[2,+\infty)$。

\end{question}
\begin{question}
\gkti{4}用反证法证明命题“设 $a,b$ 为实数，则方程 $x^3+ax+b=0$ 至少有一个实根”时，要做的假设是\mbox{（　　）}

（A）方程 $x^3+ax+b=0$ 没有实根；（B）方程 $x^3+ax+b=0$ 至多有一个实根；（C）方程 $x^3+ax+b=0$ 至多有两个实根；（D）方程 $x^3+ax+b=0$ 恰好有两个实根。

\end{question}
\begin{question}
\gkti{5}已知实数 $x,y$ 满足 $a^x<a^y$（$0<a<1$），则下列关系式恒成立的是\mbox{（　　）}

（A）$x^3>y^3$；（B）$\sin x>\sin y$；（C）$\ln(x^2+1)>\ln(y^2+1)$；（D）$\dfrac1{x^2+1}>\dfrac1{y^2+1}$。

\Needspace{0.4\textheight}
\end{question}
\begin{question}
\gkti{6}已知函数 $y=\log_a(x+c)$（$a,c$ 为常数，其中 $a>0$，$a\neq1$）的图象如图，则下列结论成立的是\mbox{（　　）}

（A）$a>1$，$c>1$；（B）$a>1$，$0<c<1$；（C）$0<a<1$，$c>1$；（D）$0<a<1$，$0<c<1$。
\nopagebreak[4]
\begin{center}% ============================================================
% f2014_shandong_arts_q06_log.tex
% 2014 山东卷（文科）第 6 题：函数 y=log_a(x+c) 的图象
% 来源 PDF：普通高考/2014/2014山东文.pdf
% 原页：第 1 页（A3 横排）左栏，印面 909，第 6 题题干下方
% 特征：图象递减，故 0<a<1；x=0 处函数值为正，故 0<c<1；
%       曲线与 x 轴交点位于 x=1 左侧，即 1-c<1。
% ============================================================
\begin{tikzpicture}[>=Stealth, line width=0.7pt, font=\small,
  x=1cm, y=1cm]

  % 坐标轴
  \draw[->] (-1.35,0) -- (4.95,0) node[below] {$x$};
  \draw[->] (0,-2.55) -- (0,3.05) node[left] {$y$};
  \node[below left] at (0,0) {$O$};

  % 曲线 y=log_{1/2}(x+0.6)，定义域 x>-0.6
  \draw[domain=-0.45:4.75, samples=220, smooth, variable=\x]
    plot ({\x}, {-ln(\x+0.6)/ln(2)});

  % x 轴上的刻度 1
  \draw (1,0.06) -- (1,-0.06);
  \node[above] at (1,0.06) {$1$};

\end{tikzpicture}
\end{center}
\nopagebreak[4]
\end{question}
\begin{question}
\gkti{7}已知向量 $\boldsymbol a=(1,\sqrt3)$，$\boldsymbol b=(3,m)$。若向量 $\boldsymbol a$，$\boldsymbol b$ 的夹角为 $\dfrac\pi6$，则实数 $m=$\mbox{（　　）}

（A）$2\sqrt3$；（B）$\sqrt3$；（C）$0$；（D）$-\sqrt3$。

\Needspace{0.4\textheight}
\end{question}
\begin{question}
\gkti{8}为了研究某种药品的疗效，选取若干名志愿者进行临床试验。所有志愿者的舒张压数据（单位：kPa）的分组区间为 $[12,13)$，$[13,14)$，$[14,15)$，$[15,16)$，$[16,17]$，将其按从左到右的顺序分别编号为第一组，第二组，…，第五组。如图所示是根据试验数据制成的频率分布直方图。已知第一组与第二组共有 $20$ 人，第三组中没有疗效的有 $6$ 人，则第三组中有疗效的人数为\mbox{（　　）}

（A）$6$；（B）$8$；（C）$12$；（D）$18$。
\nopagebreak[4]
\begin{center}% ============================================================
% f2014_shandong_arts_q08_histogram.tex
% 2014 山东卷（文科）第 8 题：舒张压数据的频率分布直方图
% 来源 PDF：普通高考/2014/2014山东文.pdf
% 原页：第 1 页（A3 横排）中栏，印面 909，第 8 题题干右侧
% 各矩形高（频率/组距）：0.24、0.16、0.36、0.16、0.08，和为 1。
% ============================================================
\begin{tikzpicture}[>=Stealth, line width=0.7pt, font=\small,
  x=1cm, y=1cm]

  % 坐标轴
  \draw[->] (0.60,0) -- (7.35,0) node[below right] {舒张压/kPa};
  \draw[->] (0.60,-0.05) -- (0.60,4.05) ;
  \node[below left] at (0.60,0) {$0$};

  % 纵轴标题：频率 / 组距（上下两行，原图无分数线）
  \node at (0.62,4.62) {频率};
  \node at (0.62,4.16) {组距};

  % x 轴断开符号
  \draw (0.86,0) -- (0.98,0.10) -- (1.10,-0.10) -- (1.22,0.10) -- (1.34,0);

  % 直方图：组距 1，映射 12→1.90，13→2.90，…，17→6.90
  \draw (1.90,0) rectangle (2.90,2.40);
  \draw (2.90,0) rectangle (3.90,1.60);
  \draw (3.90,0) rectangle (4.90,3.60);
  \draw (4.90,0) rectangle (5.90,1.60);
  \draw (5.90,0) rectangle (6.90,0.80);

  % 纵轴刻度与虚线
  \foreach \h/\y in {0.08/0.80, 0.16/1.60, 0.24/2.40, 0.36/3.60}
  {
    \draw (0.54,\y) -- (0.66,\y);
    \node[left] at (0.52,\y) {$\h$};
    \draw[dashed] (0.66,\y) -- (6.90,\y);
  }

  % 横轴刻度
  \foreach \x/\l in {1.90/12, 2.90/13, 3.90/14, 4.90/15, 5.90/16, 6.90/17}
  {
    \draw (\x,0.05) -- (\x,-0.05);
    \node[below] at (\x,-0.02) {$\l$};
  }

\end{tikzpicture}
\end{center}
\nopagebreak[4]
\end{question}
\begin{question}
\gkti{9}对于函数 $f(x)$，若存在常数 $a\neq0$，使得 $x$ 取定义域内的每一个值，都有 $f(x)=f(2a-x)$，则称 $f(x)$ 为准偶函数。下列函数中是准偶函数的是\mbox{（　　）}

（A）$f(x)=\sqrt x$；（B）$f(x)=x^3$；（C）$f(x)=\tan x$；（D）$f(x)=\cos(x+1)$。

\end{question}
\begin{question}
\gkti{10}已知 $x,y$ 满足约束条件 $\begin{cases}x-y-1\leqslant0,\\2x-y-3\geqslant0,\end{cases}$ 当目标函数 $z=ax+by$（$a>0$，$b>0$）在该约束条件下取到最小值 $2\sqrt5$ 时，$a^2+b^2$ 的最小值为\mbox{（　　）}

（A）$5$；（B）$4$；（C）$\sqrt5$；（D）$2$。

\end{question}
\gkpart{二、填空题}

\Needspace{0.4\textheight}
\begin{question}
\gkti{11}执行如图所示的程序框图，若输入的 $x$ 的值为 $1$，则输出的 $n$ 的值为\gkblank{}。
\nopagebreak[4]
\begin{center}% ============================================================
% f2014_shandong_arts_q11_flowchart.tex
% 2014 山东卷（文科）第 11 题：程序框图
% 来源 PDF：普通高考/2014/2014山东文.pdf
% 原页：第 1 页（A3 横排）中栏，印面 909，第 11 题题干下方
% 结构：开始 → 输入 x → n=0 → 判断 x^2-4x+3≤0；
%       是：x=x+1 → n=n+1 → 返回判断；否：输出 n → 结束。
% ============================================================
\begin{tikzpicture}[>=Stealth, line width=0.7pt, font=\small,
  x=1cm, y=1cm, line cap=round, line join=round]

  % 开始
  \draw[rounded corners=6pt] (-0.95,0.42) rectangle (0.95,-0.42);
  \node at (0,0) {开始};
  \draw[->] (0,-0.42) -- (0,-0.87);

  % 输入 x
  \draw (-1.45,-0.87) -- (1.45,-0.87) -- (1.25,-1.72) -- (-1.65,-1.72) -- cycle;
  \node at (0,-1.30) {输入 $x$};
  \draw[->] (0,-1.72) -- (0,-2.16);

  % n=0
  \draw (-1.10,-2.16) rectangle (1.10,-2.96);
  \node at (0,-2.56) {$n=0$};
  \draw[->] (0,-2.96) -- (0,-3.45);

  % 判断框
  \draw (0,-3.45) -- (2.55,-4.55) -- (0,-5.65) -- (-2.55,-4.55) -- cycle;
  \node at (0,-4.55) {$x^2-4x+3\leqslant 0$};

  % 是：向下进入循环体
  \draw[->] (0,-5.65) -- (0,-6.10);
  \node at (0.45,-5.90) {是};
  \draw (-1.10,-6.10) rectangle (1.10,-6.90);
  \node at (0,-6.50) {$x=x+1$};

  \draw[->] (0,-6.90) -- (0,-7.35);
  \draw (-1.10,-7.35) rectangle (1.10,-8.15);
  \node at (0,-7.75) {$n=n+1$};

  % 回边：向下、向左、向上，再水平进入判断框左顶点
  \draw[->] (-1.10,-7.75) -- (-3.45,-7.75) -- (-3.45,-4.55) -- (-2.55,-4.55);

  % 否：向右输出
  \draw[->] (2.55,-4.55) -- (3.95,-4.55);
  \node at (3.20,-4.20) {否};
  \draw (3.95,-4.10) -- (6.35,-4.10) -- (6.15,-5.00) -- (3.75,-5.00) -- cycle;
  \node at (5.05,-4.55) {输出 $n$};
  \draw[->] (5.05,-5.00) -- (5.05,-5.45);

  % 结束
  \draw[rounded corners=6pt] (4.10,-5.45) rectangle (6.00,-6.29);
  \node at (5.05,-5.87) {结束};

\end{tikzpicture}
\end{center}
\nopagebreak[4]
\end{question}
\begin{question}
\gkti{12}函数 $y=\dfrac{\sqrt3}2\sin2x+\cos^2x$ 的最小正周期为\gkblank{}。

\end{question}
\begin{question}
\gkti{13}一个六棱锥的体积为 $2\sqrt3$，其底面是边长为 $2$ 的正六边形，侧棱长都相等，则该六棱锥的侧面积为\gkblank{}。

\end{question}
\begin{question}
\gkti{14}圆心在直线 $x-2y=0$ 上的圆 $C$ 与 $y$ 轴的正半轴相切，圆 $C$ 截 $x$ 轴所得弦的长为 $2\sqrt3$，则圆 $C$ 的标准方程为\gkblank{}。

\end{question}
\begin{question}
\gkti{15}已知双曲线 $\dfrac{x^2}{a^2}-\dfrac{y^2}{b^2}=1$（$a>0$，$b>0$）的焦距为 $2c$，右顶点为 $A$，抛物线 $x^2=2py$（$p>0$）的焦点为 $F$。若双曲线截抛物线的准线所得线段长为 $2c$，且 $|FA|=c$，则双曲线的渐近线方程为\gkblank{}。

\end{question}
\gkpart{三、解答题}

\begin{question}
\gkti{16}海关对同时从 $A$，$B$，$C$ 三个不同地区进口的某种商品进行抽样检测，从各地区进口此商品的数量（单位：件）如下表所示。工作人员用分层抽样的方法从这些商品中共抽取 $6$ 件样品进行检测。

\begin{center}
\begin{tabular}{llll}
\toprule
地区 & $A$ & $B$ & $C$ \\
\midrule
数量 & $50$ & $150$ & $100$ \\
\bottomrule
\end{tabular}
\end{center}

（1） 求这 $6$ 件样品中来自 $A$，$B$，$C$ 各地区商品的数量；

（2） 若在这 $6$ 件样品中随机抽取 $2$ 件送往甲机构进行进一步检测，求这 $2$ 件商品来自相同地区的概率。

\end{question}
\begin{question}
\gkti{17}$\triangle ABC$ 中，角 $A$，$B$，$C$ 所对的边分别为 $a$，$b$，$c$。已知 $a=3$，$\cos A=\dfrac{\sqrt6}3$，$B=A+\dfrac\pi2$。

（1） 求 $b$ 的值；

（2） 求 $\triangle ABC$ 的面积。

\Needspace{0.4\textheight}
\end{question}
\begin{question}
\gkti{18}如图，四棱锥 $P-ABCD$ 中，$AP\perp$ 平面 $PCD$，$AD\parallel BC$，$AB=BC=\dfrac12AD$，$E$，$F$ 分别为线段 $AD$，$PC$ 的中点。
\nopagebreak[4]
\begin{center}% ============================================================
% f2014_shandong_arts_q18_pyramid.tex
% 2014 山东卷（文科）第 18 题：四棱锥 P-ABCD
% 来源 PDF：普通高考/2014/2014山东文.pdf
% 原页：第 2 页（A3 横排）左栏，印面 910，第 18 题题干下方
% 结构：AP⊥平面 PCD，AD∥BC，AB=BC=AD/2；
%       E 为 AD 中点，F 为 PC 中点；连 BE、BF、EF，
%       AD、AC 为不可见棱（虚线）。
% ============================================================
\begin{tikzpicture}[>=Stealth, line width=0.7pt, font=\small,
  x=1cm, y=1cm, line cap=round, line join=round]

  \coordinate (P) at (2.15,3.90);
  \coordinate (A) at (0.35,1.55);
  \coordinate (B) at (1.25,0.55);
  \coordinate (C) at (3.55,1.05);
  \coordinate (D) at (4.85,2.10);

  % 中点
  \coordinate (E) at (2.60,1.825);
  \coordinate (F) at (2.85,2.475);

  % 不可见棱
  \draw[dashed] (A) -- (D);
  \draw[dashed] (A) -- (C);
  \draw[dashed] (E) -- (F);

  % 底面与侧棱
  \draw (P) -- (A);
  \draw (P) -- (B);
  \draw (P) -- (C);
  \draw (P) -- (D);
  \draw (A) -- (B);
  \draw (B) -- (C);
  \draw (C) -- (D);

  % 辅助线
  \draw (B) -- (E);
  \draw (B) -- (F);

  % 顶点标注
  \node[above] at (P) {$P$};
  \node[left] at (A) {$A$};
  \node[below] at (B) {$B$};
  \node[below right] at (C) {$C$};
  \node[right] at (D) {$D$};
  \node[below right] at (E) {$E$};
  \node[above right] at (F) {$F$};

\end{tikzpicture}
\end{center}
\nopagebreak[4]
（1） 求证：$AP\parallel$ 平面 $BEF$；

（2） 求证：$BE\perp$ 平面 $PAC$。

\end{question}
\begin{question}
\gkti{19}在等差数列 $\{a_n\}$ 中，已知公差 $d=2$，$a_2$ 是 $a_1$ 与 $a_4$ 的等比中项。

（1） 求数列 $\{a_n\}$ 的通项公式；

（2） 设 $b_n=a_{\frac{n(n+1)}2}$，记 $T_n=-b_1+b_2-b_3+b_4-\cdots+(-1)^nb_n$，求 $T_n$。

\end{question}
\begin{question}
\gkti{20}设函数 $f(x)=a\ln x+\dfrac{x-1}{x+1}$，其中 $a$ 为常数。

（1） 若 $a=0$，求曲线 $y=f(x)$ 在点 $(1,f(1))$ 处的切线方程；

（2） 讨论函数 $f(x)$ 的单调性。

\end{question}
\begin{question}
\gkti{21}在平面直角坐标系 $xOy$ 中，椭圆 $C:\dfrac{x^2}{a^2}+\dfrac{y^2}{b^2}=1$（$a>b>0$）的离心率为 $\dfrac{\sqrt3}2$，直线 $y=x$ 被椭圆 $C$ 截得的线段长为 $\dfrac{4\sqrt{10}}5$。

（1） 求椭圆 $C$ 的方程；

（2） 过原点的直线与椭圆 $C$ 交于 $A$，$B$ 两点（$A$，$B$ 不是椭圆 $C$ 的顶点）。点 $D$ 在椭圆 $C$ 上，且 $AD\perp AB$，直线 $BD$ 与 $x$ 轴、$y$ 轴分别交于 $M$，$N$ 两点。

　　① 设直线 $BD$，$AM$ 的斜率分别为 $k_1$，$k_2$，证明存在常数 $\lambda$ 使得 $k_1=\lambda k_2$，并求出 $\lambda$ 的值；

　　② 求 $\triangle OMN$ 面积的最大值。

\end{question}